\chapter{总结与展望}

\section{工作总结}

本文围绕实时约束下的可部署视频编码优化这一目标，研究了强化学习前处理与NVENC硬件编码器的协同机制。主要工作总结如下：

\begin{enumerate}
  \item 提出了一套包含前处理与后处理概念位的混合编码框架。本文重点对NVENC硬件编码前处理模块进行了实证验证，将不可导的闭源硬件编码器纳入优化回路，实现了前处理模型对编码器特性的感知与预补偿；

  \item 引入了基于Actor-Critic架构的RAPNet前处理网络，采用残差学习策略将动作空间约束为对原始帧的微调修正，并设计了基于相对率失真增益的复合奖励函数，使策略能够在质量与码率之间进行内容相关的自适应权衡；

  \item 提出了包含加权MSE损失、非边缘高频抑制损失和全变分平滑损失的频域感知复合损失策略，用于抑制增强过程中的高频伪影并提升输出稳定性；

  \item 在5个HEVC标准测试序列上，本文方法取得平均$-9.48\%$的BD-LPIPS收益，其中ParkScene和BasketballDrive分别达到$-13.45\%$与$-12.84\%$。在复杂度方面，模型参数量为139K、推理时延为5.4ms/帧，前处理开销约占端到端时延的12\%。这意味着在本文采用的NVENC实时编码链路中，额外引入5.4ms/帧前处理后，能够换取可测的感知率失真收益。
\end{enumerate}

\section{研究局限性与讨论}

在本文的研究与评测过程中，笔者也充分认识到了当前方案在理论与工程实践上的部分局限性，在此作补充讨论，以期为后续研究提供更严谨的视角：

\begin{enumerate}
  \item \textbf{强化学习与监督学习的权衡}：研究初期曾尝试更为直接的监督学习方案（如RPP），但后续研究转向强化学习。这一决策的根本原因在于应用场景，闭源NVENC硬件编码器不可导，且其码率控制算法对平滑信号极其敏感。监督学习缺乏来自真实编码器的环境反馈，容易出现客观指标较高但不适配硬件RC的情况。强化学习虽训练成本高、收敛更难，但在无梯度环境下通过实测反馈进行探索，在本文任务中更适合作为前处理与硬件编码器协同优化的方法。
  \item \textbf{客观评价指标的选取博弈}：本文在主要结果评测中着重汇报了LPIPS / BD-LPIPS，而非传统的PSNR。这是因为在低码率（200Kbps$\sim$800Kbps）的极端压缩工况下，重建画面通常伴随严重的块效应和振铃伪影。传统基于像素误差的PSNR往往对平滑但丧失纹理的画面给出高分，与人眼主观感知严重背离；而基于深度特征提取的LPIPS对结构失真和高频伪影有着更准确的惩罚机制，更能反映真实观感。
  \item \textbf{数据规模与外推的局限性}：当前系统仅使用5个UVG标准测试序列（1080p）进行了定量验证。虽然这类测试序列在课题研究中具有代表性，且覆盖了平坦、快速运动、复杂纹理等典型场景，但这毕竟属于中小规模评测。当前结论主要成立在针对HEVC编码、1080p分辨率与特定帧率的实时工程场景下。若要将结论外推至4K/8K超高清、非自然生成视频（如游戏录屏）或其他硬件平台，仍需更大规模的数据集与多样化编码器的进一步实验验证。
\end{enumerate}

\section{未来工作展望}

尽管本文已在既定实验条件下验证了方案的有效性，但受限于本科阶段可获取的计算资源、可覆盖的数据规模以及异构硬件平台条件，仍有若干更具系统性的工作尚未充分展开。后续研究可重点围绕以下三个方向推进：

\begin{enumerate}
  \item \textbf{更大规模与多分辨率的数据验证}：当前实验主要基于1080p标准测试序列展开，能够支持方法有效性的初步论证，但对更广泛内容分布的覆盖仍然有限。若要进一步扩展到2K、4K乃至更复杂场景，需要构建规模更大的训练与测试集，并完成大量离线编码与指标统计。这一过程对GPU算力、存储空间和实验周期都有较高要求，超出了当前毕业设计阶段的可承受范围，因此本文暂未进行更大规模的数据扩充与系统验证；

  \item \textbf{跨编码器与跨硬件平台的泛化评估}：本文工作主要围绕NVENC实时编码链路展开，结论目前仍建立在特定硬件平台与码率控制机制之上。若要将方法进一步推广到VVC、AV1或不同厂商的硬件编码器，除需重新适配实验链路外，还需要长期稳定获取多种GPU、移动端SoC或专用编解码设备。受当前平台条件限制，本文尚不具备完成大范围异构平台对比实验的硬件基础，后续可在条件允许时系统评估该方法的跨平台迁移能力；

  \item \textbf{面向真实系统的端到端部署研究}：离线实验已经表明前处理模块能够带来稳定的感知收益，但若要进入实际视频通信或流媒体系统，还需进一步考察端到端时延、显存占用、功耗波动与长期运行稳定性等工程指标。这类研究不仅需要完整的业务系统支撑，还依赖持续的联调、部署和性能剖析条件，而这些工作量已明显超出本文以算法验证为主的研究边界。后续若具备更完整的工程环境，可进一步开展真实业务场景下的部署验证与联合优化。
\end{enumerate}
